\documentclass[11pt]{article}
\usepackage[a4paper,margin=25mm]{geometry}
\usepackage{amsmath,amssymb,amsthm}
\usepackage[hidelinks]{hyperref}
\emergencystretch=2em
\newtheorem{theorem}{Theorem}
\title{Sector angles do not determine a calculus bound\newline
\large Disproof of Conjecture 00000002362}
\author{gaochengzhecpu}\date{}
\begin{document}\maketitle
\begin{abstract}
On the fixed Hilbert space $\mathbb C^2$ with its Euclidean
norm, consider $A_N=I+NE_{12}$, $N\geq0$. All these
operators are sectorial of angle zero, hence belong to
the same positive-angle sectorial class. On the fixed
larger sector given by the right half-plane, the bounded
holomorphic rational function $g(z)=z/(1+z)^3$ has
supremum norm at most one, but
$g(A_N)=I/8-NE_{12}/16$ has unbounded operator norm.
Thus no finite calculus bound can depend only on the
sector angles. Each individual operator does have a
bounded calculus; the obstruction is uniformity across
operators with the same angles.
\end{abstract}

\section*{A fixed Hilbert space and fixed sector angles}
Write $H=\mathbb C^2$ with its Euclidean norm, and use
the genuine operator norm on $\mathcal B(H)$. Set
\[
 E=\begin{pmatrix}0&1\\0&0\end{pmatrix},\qquad
 A_N=I+NE,\qquad N\geq0.
\]
Then $E^2=0$ and $\|E\|=1$. For $0\leq\omega<\pi$,
the standard sectorial condition is spectral containment
in $\overline\Sigma_\omega$ together with a finite bound
\[
 \sup_{z\notin\overline\Sigma_\theta}
                \|z(zI-A)^{-1}\|<\infty
                 \quad(\omega<\theta<\pi),
\]
where $\Sigma_\theta=\{z\ne0:|\arg z|<\theta\}$.
This is spectral/resolvent sectoriality, not a numerical
range condition; see \cite[Section 2.1]{bhm}.

We prove $A_N$ is sectorial of angle zero for every $N$.
In particular all $A_N$ lie in the same class
$\operatorname{Sect}(\pi/3)$, and we take their calculus
on the fixed strictly larger sector
$\Sigma_{\pi/2}=\{z:\operatorname{Re}z>0\}$.
The class angle $\pi/3$ is not claimed to be their
least angle. Their stronger angle-zero membership
also prevents an operator-dependent choice of angle
from explaining the growing constants.

\newpage
\section*{The genuine spectrum and resolvent condition}
For every $z\ne1$, the actual resolvent is
\[
 (zI-A_N)^{-1}=(z-1)^{-1}I+N(z-1)^{-2}E.
\]
Multiplying by $zI-A_N$ verifies the identity, using
$E^2=0$. Thus $\sigma(A_N)\subseteq\{1\}$.
In fact $A_N-I=NE$ is not invertible, so its spectrum
is exactly $\{1\}$; the inclusion alone suffices below.

Fix $0<\theta<\pi$ and let
$\eta=\min(1,1-\cos\theta)>0$. If
$z\notin\overline\Sigma_\theta$, then
$\operatorname{Re}z\leq|z|\cos\theta$. Therefore
\[
 |z-1|^2=|z|^2+1-2\operatorname{Re}z,
                         \qquad |z-1|\geq\eta.
\]
To see the second inequality, put $r=|z|$ and
$c=\cos\theta$. If $c\geq0$, the square is at least
$(r-c)^2+1-c^2\geq(1-c)^2$; the last comparison uses
$2c(1-c)\geq0$. If $c<0$, then
$\operatorname{Re}z\leq0$, so the square is at least one.

Since $|z|\leq|z-1|+1$ and $\|E\|=1$, the resolvent
identity gives the genuine uniform-in-$z$ estimate
\begin{align*}
 \|z(zI-A_N)^{-1}\|
  &\leq\frac{|z|}{|z-1|}
           +N\frac{|z|}{|z-1|^2}\\
  &\leq(1+\eta^{-1})(1+N\eta^{-1}).
\end{align*}
For each $N$ this is finite for every larger sector.
Together with spectral containment on the positive
axis, this proves sectoriality of angle zero.
The resolvent constants are permitted to depend on
$N$ in this definition; that dependence is precisely
what angle information alone fails to control.

\section*{An admissible bounded holomorphic rational test}
Take $g(z)=z/(1+z)^3$. Its only pole is at $-1$, so it
is holomorphic throughout the right half-plane.
For $\operatorname{Re}z\geq0$,
\[
 |1+z|\geq1,\qquad
 |1+z|^2\geq1+|z|^2.
\]
Consequently
\[
 |g(z)|\leq\frac{|z|}{1+|z|^2}\leq1,
 \qquad
 |g(z)|\leq\frac{2|z|}{(1+|z|)^2}.
\]
Thus $g$ is also in the usual decaying elementary
class $H_0^\infty$: it vanishes suitably at zero and
infinity. No issue about extending the calculus to
constants or nondecaying functions is needed.

\newpage
\section*{The rational calculus has unbounded norm}
The actual inverse of $I+A_N=2I+NE$ is
\[
               (I+A_N)^{-1}=\tfrac12I-\tfrac N4E.
\]
Using $E^2=0$ once more gives its ordinary rational
functional-calculus value
\[
             g(A_N)=A_N(I+A_N)^{-3}
                         =\tfrac18I-\tfrac N{16}E.
\]
The triangle inequality and $\|E\|=\|I\|=1$ imply
\[
                    \|g(A_N)\|\geq\tfrac N{16}-\tfrac18.
\]
For any proposed finite real constant $C$, choose
$N=16(\max(C,0)+1)$. The right side strictly exceeds
$C$, whereas $\|g\|_{\infty,\Sigma_{\pi/2}}\leq1$.

Any bounded $H^\infty$ calculus agrees with the ordinary
rational calculus and obeys
\[
       \|f(A)\|\leq K\|f\|_{\infty,\Sigma_{\pi/2}}.
\]
Hence its bound $K$ must dominate $\|g(A_N)\|$ for
this single admissible test. No finite constant
determined only by the fixed angles $\pi/3$ and
$\pi/2$ can do so for all $N$.

\section*{Each operator individually has a bounded calculus}
This example does not depend on choosing operators
without a bounded calculus. For each fixed $N$, the
ordinary holomorphic calculus at this Jordan block is
\[
                 f(A_N)=f(1)I+Nf'(1)E.
\]
Indeed, substituting the above resolvent into the
holomorphic contour formula gives the value $f(1)$
from the simple pole and $f'(1)$ from the double pole.
The circle $|z-1|=1/2$ lies in the right half-plane.
Cauchy's derivative estimate therefore gives
$|f'(1)|\leq2\|f\|_\infty$, and consequently
\[
                  \|f(A_N)\|\leq(1+2N)\|f\|_\infty.
\]
Thus each has a legitimate finite bound, but the bounds
cannot be uniform over this fixed-angle family.
This last contour/Cauchy explanation is the mathematical
bridge to the full calculus; the Lean project checks
the necessary rational-test obstruction directly.

\begin{theorem}
Even on the fixed Hilbert space $\mathbb C^2$, no finite
bound for the sectorial $H^\infty$ calculus can depend
only on the sectorial and calculus-domain angles.
\end{theorem}
\begin{proof}
The family $A_N$ and the single test $g$ above have
fixed angles and bounded scalar norm but unbounded
operator values, contradicting every such bound.
\end{proof}

\newpage
\section*{Scope, formalization and reproducibility}
This refutes the angle-only constant assertion in the
source, already at fixed angles. There is no need to
analyze its additional claimed singularity as the angle
tends to $\pi$. Estimates involving resolvent bounds,
similarity constants or additional operator hypotheses
are not refuted.

The self-contained Lean project first works in complex
$2\times2$ matrices with Mathlib's actual Euclidean
operator norm. It proves $E^2=0$, $\|E\|=1$, both
displayed inverses, the true spectrum containment,
the all-larger-sectors resolvent estimates, and the
rational value and norm lower bound. It also proves
the actual half-plane/sector equality, scalar holomorphy,
uniform scalar bound and both displayed decay estimates.

Mathlib's star algebra equivalence from matrices to
continuous complex-linear operators on Euclidean
$\mathbb C^2$ preserves the norm. The project explicitly
transfers spectrum, units, resolvents, sectoriality and
the rational expression to this genuine operator
algebra. Invertibility of $I+A_N$ is proved; a totalized
inverse of a singular operator is never used.

The final formal statement says that for every real
$C$ there is an actual operator in
$\operatorname{Sect}(\pi/3)$ whose rational-test norm
exceeds $C$. A separate theorem negates the existence
of a uniform finite bound. The scalar function is
formally shown admissible on the same fixed larger
domain. These prove a necessary consequence of the
claimed full-calculus estimate false; the project
does not redefine a substitute $H^\infty$ calculus or
claim to formalize its entire contour construction.

Lean 4.19.0 and Mathlib revision\\
\texttt{c44e0c8ee63ca166450922a373c7409c5d26b00b}
are pinned. From \texttt{lean/}, run
\begin{verbatim}
lake exe cache get
lake build
lake env lean -DwarningAsError=true Main.lean
\end{verbatim}
Actual fresh-build logs, exact source and hashes are
included. Printed dependencies contain only standard
foundational axioms; no auxiliary numerical verifier
is needed. Compile with \texttt{tectonic main.tex}.
The author agent and parent reviewed the work internally;
no external independent review is claimed.

\begin{thebibliography}{9}
\bibitem{source} The Last Math Competition,
Conjecture 00000002362. Exact bilingual source included.
\bibitem{bhm} C. Batty, M. Haase and J. Mubeen,
\emph{The Holomorphic Functional Calculus Approach to
Operator Semigroups}, September 20, 2012.
\href{https://fa.ewi.tudelft.nl/~haase/files/bahamu20-09-2012.pdf}{Author-hosted paper}.
Section 2.1 gives sectoriality; Section 5 specifies
the calculus norm estimate. The construction above
supplies its own counterexample and estimates.
\end{thebibliography}
\end{document}
